% Generated draft: requires mathematical and source review.
\documentclass[11pt]{article}
\usepackage[margin=25mm]{geometry}
\usepackage{fontspec}
\setmainfont{Latin Modern Roman}
\usepackage{amsmath,amssymb,mathtools,unicode-math}
\setmathfont{Latin Modern Math}
\usepackage{enumitem,needspace,xurl,hyperref}
\setlength{\parindent}{0pt}
\setlength{\parskip}{5pt}
\setlength{\emergencystretch}{3em}
\newcommand{\N}{\mathbb{N}}
\newcommand{\Z}{\mathbb{Z}}
\newcommand{\Q}{\mathbb{Q}}
\newcommand{\R}{\mathbb{R}}
\newcommand{\C}{\mathbb{C}}
\newcommand{\F}{\mathbb{F}}
\newcommand{\A}{\mathbb{A}}
\newcommand{\PP}{\mathbb{P}}
\newcommand{\E}{\mathbb{E}}
\newcommand{\eps}{\varepsilon}
\newcommand{\dd}{\,\mathrm{d}}
\begin{document}
\section*{UnsolvedMath ID: 20001662}
\textbf{Source problem code:} catalogue entry
\subsection{Definitions and mathematical statement}
For a finitely generated discrete group \(G\), choose a finite symmetric generating set and word length \(L\). Property RD means that there are \(C,s>0\) such that, for every finitely supported \(f\) with support in \(\{g:L(g)\leq R\}\),
\[
 \|\lambda_G(f)\|\leq C(1+R)^s\|f\|_2,
\]
where \(\lambda_G(f)\) is the left convolution operator on \(\ell^2(G)\). The recovered Section 8 question asks which of these groups have RD: \(\operatorname{Out}(F_n)\), \(\operatorname{Aut}(F_n)\), braid groups \(B_n\), mapping class groups, and Artin groups. The original note records \(B_3\), right-angled Artin groups and Coxeter groups as positive cases. Its brief cubical explanation for RAAGs must not be read as a general assertion that every free action on an arbitrary cube complex implies RD.

The aid reports that Behrstock--Minsky settled the finite-type orientable mapping-class-group and all-braid-group items affirmatively. It describes \(\operatorname{Aut}(F_n)\) and \(\operatorname{Out}(F_n)\) as known through rank two and apparently open for \(n\geq3\), and the general Artin question as open despite many positive subclasses. These are attributed literature assessments, with no complete references supplied.

The aid's quantitative contribution concerns induced cyclic length. Let \(z\in Z(G)\) have infinite order, and let \(\varepsilon:G\to\Z\) be a homomorphism with \(\varepsilon(z)\ne0\). Put \(M=\max_{s\in S}|\varepsilon(s)|>0\). On \(\langle z\rangle\) use the restriction of the ambient word length, rather than silently substituting its intrinsic cyclic word length. The reported estimates are
\[
 L(z^k)\geq\frac{|k\varepsilon(z)|}{M},\qquad
 \#\{z^k:L(z^k)\leq R\}
 \leq2\left\lfloor\frac{MR}{|\varepsilon(z)|}\right\rfloor+1.
\]
For a function \(f\) on \(\langle z\rangle\) supported in that induced ball, it gives
\[
 \|\lambda_{\langle z\rangle}(f)\|
 \leq\left(2\left\lfloor\frac{MR}{|\varepsilon(z)|}\right\rfloor+1\right)^{1/2}\|f\|_2.
\]
This verifies RD for the induced kernel length, with an explicit polynomial bound, not merely RD for an unrelated intrinsic length.

For \(B_n\), \(n\geq2\), use the standard generators \(\sigma_1,\ldots,\sigma_{n-1}\), exponent sum \(\varepsilon(\sigma_i)=1\), and full twist
\[
 z_n=\Delta^2=(\sigma_1\cdots\sigma_{n-1})^n.
\]
The aid reports the sharp formula
\[
 L(z_n^k)=n(n-1)|k|,\qquad
 \#\{z_n^k:L(z_n^k)\leq R\}
 =2\left\lfloor\frac{R}{n(n-1)}\right\rfloor+1.
\]
With this induced-length certificate and RD of the quotient in
\[
 1\to\langle z_n\rangle\to B_n\to B_n/\langle z_n\rangle\to1,
\]
the aid invokes Garncarek's extension theorem to obtain a separate central-extension derivation of RD for every braid group. For \(n=1\), the braid group is trivial and the infinite-central-element formula is not applicable. No complete extension theorem or quotient identification is supplied in the summary.
\subsection{Short English statement}
The recovered question asks RD for five families, rather than merely listing groups. The aid reports the mapping-class and braid cases solved and supplies an exact induced-length calculation for braid full twists, sufficient for a cited central-extension argument. The outer/free automorphism and general Artin statuses remain attributed open assessments.
\subsection{Sources}
\begin{itemize}
\item Local UnsolvedMath registry, permanent ID \(20001662\), code catalogue entry; canonical \texttt{aim-geometric-group-theory-notes.json}, \texttt{notes[170]}.
\item AIM workshop notes ``The property of rapid decay,'' 21 March 2006, Section 8, \url{https://aimath.org/WWN/rapiddecay/rapiddecay.pdf}. The supplied verification locates the heading ``Which one of these groups have RD?'' at the preceding page boundary, absorbed into corpus record0170 as ``68,'' combining page number6 and section number8. No PDF inspection is claimed here.
\item Behrstock--Minsky are named for RD of mapping class and braid groups; Garncarek is named for the extension criterion. No paper titles, dates, theorem numbers or complete authoritative references are supplied.
\item Research attempt 1: ``partial\_result,'' ``partially\_solved,'' rigor label ``proved,'' novelty confidence ``low.'' The aid is machine generated, not peer reviewed; independently verify before citation.
\end{itemize}
\end{document}
